
We investigated whether preference entropy could serve as an information-theoretic proxy for bidirectional alignment in RLHF evaluation---measuring not only whether AI aligns to humans (preference agreement) but whether humans preserve critical evaluation capacity when interacting with aligned AI (preference diversity). Our validation reveals a methodological limitation: standard datasets (Anthropic-HH, WebGPT) use pairwise comparison formats optimized for annotation cost, structurally incompatible with per-prompt entropy measurement.

\textbf{Key Findings:}

\begin{enumerate}
\item \textbf{Dataset Structure Incompatibility:} Pairwise-unique response formats enable efficient reward model training but prevent diversity measurement. Entropy aggregation yields tautological 50/50 splits (constant H = ln(2) $\approx$ 0.693 nats), not actual preference distributions.
\end{enumerate}

\begin{enumerate}
\item \textbf{Technical Validation Success:} Entropy computation mechanism validated (100\% success rate, correct mathematical range [0, ln(2)]), proving feasibility with appropriate data structure. Issue is dataset format compatibility, not implementation capability.
\end{enumerate}

\begin{enumerate}
\item \textbf{Dataset Format Documentation:} Documented cost-measurement trade-off between pairwise-unique (cost-efficient, entropy-incompatible) and multi-annotator-fixed (entropy-measurable, higher cost) formats. Tested on Anthropic-HH (n=100 prompts) with format-based inference for other datasets.
\end{enumerate}

\begin{enumerate}
\item \textbf{Alternative Proxies:} Proposed three measurement approaches when multi-annotator data unavailable: response diversity entropy, intra-annotator variance, entropy-regularized RLHF. All untested.
\end{enumerate}

\textbf{Contributions:} We introduce bidirectional alignment as evaluation framework, distinguish unidirectional metrics ($AI\rightarrow human$ preference agreement) from bidirectional metrics ($human\rightarrow AI$ agency preservation), document the cost-measurement trade-off in dataset structure for preference diversity measurement (tested on n=1 dataset with inference for others), and propose actionable alternatives (response diversity proxy, multi-annotator evaluation subsets, entropy-regularized training).

\textbf{Implications:} Current RLHF benchmarks cannot distinguish legitimate consensus (users agree because AI is correct on objective tasks) from preference homogenization (users agree because they stopped critically evaluating on subjective tasks). Without diversity metrics, we are blind to potential bidirectional alignment failures where RLHF succeeds at making users agree but fails at preserving agency on tasks where diverse preferences are legitimate.

\textbf{Future Work:} Three research directions emerge. Short-term (2-4 weeks): Validate response diversity entropy as entropy collapse proxy. Medium-term (2-3 months): Collect multi-annotator preference data (1K prompts $\times$ 20 annotators) to re-validate H-E1 with appropriate dataset structure. Long-term (6-12 months): Develop and evaluate entropy-regularized RLHF algorithms that preserve diversity on subjective tasks while allowing convergence on objective tasks.

The hypothesis mechanism (entropy collapse beyond performance-justified levels on subjective tasks) remains untested due to dataset limitation, not falsified. Future work with appropriate data structures can test whether RLHF inadvertently imposes preference monoculture on subjective tasks, and whether entropy-regularized training can prevent it while preserving alignment quality.

High preference agreement could indicate successful alignment or problematic homogenization---current benchmarks cannot distinguish these scenarios. Our work identifies why (dataset structure incompatibility) and proposes how to address it (multi-annotator subsets, response diversity proxy, entropy-regularized RLHF). Distinguishing legitimate consensus from homogenization requires diversity metrics. Current benchmarks structurally cannot provide them---but they could, with deliberate design choices prioritizing bidirectional alignment evaluation.

